\documentclass{article}

% A synthetic seed paper for Lattice's research environment. Its numbers are
% illustrative, not measured results.
\usepackage[preprint]{icml2026}
\usepackage[utf8]{inputenc}
\usepackage[T1]{fontenc}
\usepackage{times}
\usepackage{amsmath,amssymb,amsfonts}
\usepackage{booktabs}
\usepackage{graphicx}
\usepackage{subcaption}
\usepackage{siunitx}
\usepackage{enumitem}
\usepackage{nicefrac}
\usepackage{hyperref}
\usepackage{microtype}
\usepackage{xcolor}

\newcommand{\ece}{\mathrm{ECE}}
\newcommand{\prob}[1]{\Pr\left[#1\right]}

\icmltitlerunning{Calibrated Uncertainty for Tabular Benchmarks}

\begin{document}

\twocolumn[
  \icmltitle{Calibrated Uncertainty for Tabular Benchmarks}

  \begin{icmlauthorlist}
    \icmlauthor{Seed Author}{lattice}
  \end{icmlauthorlist}

  \icmlaffiliation{lattice}{Lattice Research Environment}
  \icmlcorrespondingauthor{Seed Author}{seed@lattice.example}
  \icmlkeywords{Calibration, Tabular Data, Ensembles}

  \vskip 0.3in
]

\printAffiliationsAndNotice{}

\begin{abstract}
Deep networks trained on tabular data are often overconfident, and benchmarks that report only accuracy hide the problem.
We measure calibration on a suite of tabular classification tasks, compare temperature scaling with deep ensembles, and find that a short list of evaluation practices changes which method looks best.
\end{abstract}

\section{Introduction}
\label{sec:introduction}

A classifier is calibrated when its confidence matches its accuracy: among the inputs it labels with confidence $0.8$, about $80\%$ should be correct.
Modern deep networks are frequently miscalibrated \citep{guo2017calibration}, and on tabular data they are not always more accurate than tree ensembles either \citep{grinsztajn2022tree}.
This note asks how much the choice of calibration metric and evaluation protocol matters when comparing uncertainty methods on such data.
\autoref{sec:setup} describes the protocol, \autoref{sec:results} the results, and \autoref{sec:practices} the practices we recommend.

\section{Setup}
\label{sec:setup}

We use the expected calibration error with $B$ equal-width bins \citep{naeini2015obtaining},
\[
  \ece = \sum_{b=1}^{B} \frac{|S_b|}{N} \left| \mathrm{acc}(S_b) - \mathrm{conf}(S_b) \right|,
\]
where $S_b$ is the set of test points whose confidence falls into bin $b$.
A perfectly calibrated model satisfies $\prob{Y = \hat{Y} \mid \hat{P} = p} = p$ for every $p$.
We compare three methods:
\begin{enumerate}[label=(\alph*)]
  \item a single multilayer perceptron (MLP);
  \item the same MLP after temperature scaling on a held-out split \citep{guo2017calibration};
  \item a deep ensemble of five MLPs \citep{lakshminarayanan2017simple}.
\end{enumerate}
Each dataset is split into \num{60}/\num{20}/\num{20} percent for training, calibration and testing.
Training one MLP takes about \SI{3}{\minute} on a single CPU core.

\section{Results}
\label{sec:results}

\autoref{tab:ece} reports test accuracy and $\ece$ with $B = 15$.\footnote{All numbers in this seed paper are illustrative and were not measured.}
Temperature scaling leaves accuracy unchanged by construction and removes most of the single model's calibration error.
The ensemble is both more accurate and better calibrated, at five times the training cost.

\begin{table}[t]
  \centering
  \caption{Accuracy and calibration error on the tabular suite (illustrative).}
  \label{tab:ece}
  \begin{tabular}{@{}lcc@{}}
    \toprule
    Method & Accuracy (\%) & $\ece$ (\%) \\
    \midrule
    MLP & 84.1 & 7.9 \\
    MLP + temperature & 84.1 & 2.6 \\
    Ensemble of five & 85.7 & 1.9 \\
    \midrule
    \multicolumn{3}{@{}l}{\emph{Gradient-boosted trees, for reference}} \\
    Boosted trees & 86.3 & 3.4 \\
    \bottomrule
  \end{tabular}
\end{table}

\begin{figure}[t]
  \centering
  \begin{subfigure}[t]{0.48\linewidth}
    \centering
    \includegraphics[width=\linewidth]{figures/reliability-mlp.png}
    \caption{Single MLP.}
    \label{fig:reliability-mlp}
  \end{subfigure}\hfill
  \begin{subfigure}[t]{0.48\linewidth}
    \centering
    \includegraphics[width=\linewidth]{figures/reliability-ensemble.png}
    \caption{Ensemble of five.}
    \label{fig:reliability-ensemble}
  \end{subfigure}
  \caption{Reliability diagrams on the largest dataset. Bars below the diagonal mean overconfidence.}
  \label{fig:reliability}
\end{figure}

\autoref{fig:reliability} shows reliability diagrams for the largest dataset.
The single model is overconfident in its top bins (\autoref{fig:reliability-mlp}), while the ensemble tracks the diagonal closely (\autoref{fig:reliability-ensemble}).
Reducing the bin count to $B = 5$ lowers every method's $\ece$ by roughly $\nicefrac{1}{3}$, so results reported with different $B$ are not comparable.

\section{Recommended practices}
\label{sec:practices}

\begin{description}[style=nextline]
  \item[Report the bin count.] The value of $\ece$ depends on $B$; state it, and prefer a fixed $B$ across a paper.
  \item[Calibrate on held-out data.] Fitting the temperature on the test split overstates its benefit.
  \item[Compare at equal cost.] An ensemble of five should be compared with a single model trained five times longer, or with five calibrated single models.
\end{description}

\section{Conclusion}
\label{sec:conclusion}

On the tabular suite, temperature scaling closes most of the calibration gap at no cost to accuracy, and ensembles close the rest at a higher training cost.
Which method looks best depends on the bin count and on whether costs are matched, so both belong in every comparison.

\bibliography{references}
\bibliographystyle{icml2026}

\end{document}
